\documentclass{article}
\usepackage{geometry}
\usepackage{graphicx} % Required for inserting images
\usepackage{amsmath, amsthm, amssymb}
\usepackage{parskip}
\newgeometry{vmargin={15mm}, hmargin={24mm,34mm}}
\theoremstyle{definition} 
\newtheorem{definition}{Definition}

\newtheorem{theorem}{Theorem}[section]
\newtheorem{lemma}[theorem]{Lemma}
\newtheorem{proposition}[theorem]{Proposition}
\newtheorem{corollary}{Corollary}[theorem]


\title{MATH210B Homework 6}
\date{March 2025}
\author{Boran Erol}

\begin{document}

\maketitle

\section{Problem 8}

Let $ A : V \to V $ be a linear operator. Let $ \beta $ be a basis such that
$ [A]_{\beta} $ is diagonal.

Then, $ A(v) = \lambda v $ for any $ v \in \beta $.

If this is the case, then, $ V $ is the direct sum of eigenspaces since
eigenspaces with different eigenvalues don't intersect and since the eigenbasis generates $ V $,
they cover all of $ V $.

Then, all elementary divisors of $ A $ are linear since the characteristic polynomial splits.

This immediately implies that all invariant factors of $ A $ are products of distinct linear polynomials
since they are products of elementary divisors. Since the minimal polynomial is an invariant factor,
this also immediately implies that the minimal polynomial is a product of linear polynomials.

Lastly, notice that cyclic modules $ F[x] / (x - \lambda_{i}) $ correspond to subspaces where $ A $ acts by
multiplication by $ \lambda_{i} $, so if we can decompose $ V $ into such subspaces, then $ A $ is diagonalizable.


\section{Problem 9}

If $ A^{N} = 0 $ for some $ N > 0 $, $ x^{N} $ acts as zero. Then,
the minimal polynomial divides $ x^{N} $ and is thus $ x^{k} $ for some
$ k \leq n $, which produces the desired result.
 
\section{Problem 10}

Notice that $ det(A - tI) = det((A - tI)^{t}) = det(A^{t} - tI) $, so both matrices
have the same characteristic polynomial. But then they have the same elementary divisors
and they're thus similar.



\end{document}
